\textbf{COLUNI 2019 -- Questão 06}\par\medskip A identidade do homem/poeta e sua cultura regional estão interligadas na obra de Manoel de Barros. Assinale o excerto em que essa característica NÃO ocorre:\begin{enumerate}
\item “Na aba do campo a namorada do Sebastião torcia: quebra esse Saba, destina eles pras piranhas.”
      (BARROS, M. \textbf{Memórias inventadas}. Rio de Janeiro: Alfaguara, 2018, p. 52.)
\item “Nos seus olhos nublados eu vejo a flora dos corixos.” (BARROS, M. \textbf{Meu quintal é maior do que o mundo}. Rio de Janeiro:
      Alfaguara, 2015, p. 66)
\item “Ai morena, não me escreve / que eu não sei a ler. Aquele a preposto ao verbo ler, ao meu ouvir, ampliava
      a solidão do vaqueiro.” (BARROS, M. \textbf{Memórias inventadas}. Rio de Janeiro: Alfaguara, 2018, p. 24)
\item “Tomei um vidro de fortificante e fiquei bom. A esse tempo também eu aprendi a escutar o silêncio das
      paredes.” (BARROS, M. \textbf{Memórias inventadas}. Rio de Janeiro: Alfaguara, 2018, p. 20)\end{enumerate}
